\documentclass[a4paper,12pt,titlepage,finall]{article}

\usepackage[T1,T2A]{fontenc}
\usepackage[utf8x]{inputenc}
\usepackage[russian]{babel}
\usepackage{tikz}
\usepackage{pgfplots}
\usepackage{geometry}
\usepackage{indentfirst}
\usepackage{listings}
\usepackage{xcolor}

\geometry{a4paper,left=30mm,top=30mm,bottom=30mm,right=30mm}
\setcounter{secnumdepth}{0}
\pgfplotsset{compat=1.18}
\usepgfplotslibrary{fillbetween}

\lstset{
    basicstyle=\ttfamily\small,
    language=C,
    frame=single,
    breaklines=true,
    keepspaces=true,
    tabsize=4
}

\begin{document}

% ------------------------------------------------------------------ Титульный лист
\begin{titlepage}
    \begin{center}
        {\small \sc Московский государственный университет \\имени М.~В.~Ломоносова\\
        Факультет вычислительной математики и кибернетики\\}
        \vfill
        {\Large \sc Отчёт по заданию №6}\\[4mm]
        {\large \bf <<Сборка многомодульных программ.\\
        Вычисление корней уравнений и определённых интегралов.>>}\\[4mm]
        {\large \bf Вариант 10 / 5 / 3}
    \end{center}
    \begin{flushright}
        \vfill
        {Выполнил:\\
        студент 104 группы\\
        Москвин~Г.~Е.\\[4mm]
        Преподаватель:\\
        (указать)}
    \end{flushright}
    \begin{center}
        \vfill
        {\small Москва\\2026}
    \end{center}
\end{titlepage}

\tableofcontents
\newpage

% ================================================================== 1
\section{Постановка задачи}

Требуется разработать программу, которая вычисляет площадь плоской фигуры,
ограниченной тремя кривыми
\[
    y = f_1(x) = 1 + \frac{4}{x^2+1},\qquad
    y = f_2(x) = x^3,\qquad
    y = f_3(x) = 2^{-x},
\]
с заданной точностью $\varepsilon \in \{10^{-5},\;10^{-7},\;10^{-9}\}$.

Программа решает задачу в два этапа.
Сначала с точностью $\varepsilon_1$ численно находятся абсциссы трёх
точек попарного пересечения кривых~--- вершины фигуры.
Используются два метода нахождения корней уравнения $F(x)=0$:
\textbf{метод Ньютона (касательных)} и
\textbf{комбинированный метод (хорды + касательные)}.
Отрезки, на которых применяются эти методы, определяются
аналитически на основе анализа поведения функций.
Выбор метода осуществляется на этапе компиляции с помощью ключа
препроцессора \texttt{-DMETHOD\_NEWTON} или \texttt{-DMETHOD\_COMBINED}.

Затем площадь представляется как алгебраическая сумма определённых
интегралов, которые вычисляются с точностью $\varepsilon_2$
\textbf{методом Симпсона} с контролем погрешности по правилу Рунге.

Функции $f_1$, $f_2$, $f_3$ и их производные реализованы на ассемблере
NASM32 (архитектура i386, сопроцессор x87, соглашение~\texttt{cdecl}).
Остальная часть программы написана на языке~Си.

\newpage

% ================================================================== 2
\section{Математическое обоснование}

\subsection{Графики функций и визуализация фигуры}

На рис.~\ref{plot1} приведены графики трёх заданных кривых.
Фигура имеет три вершины~--- точки попарного пересечения.

\begin{figure}[h]
\centering
\begin{tikzpicture}
\begin{axis}[
    axis lines=middle,
    restrict x to domain=-2.5:2.5,
    restrict y to domain=-0.5:6,
    xmin=-2.5, xmax=2.5, ymin=-0.5, ymax=6,
    enlargelimits=false,
    legend cell align=left,
    legend pos=north east,
    grid=major,
    scale=1.6,
    xlabel={$x$}, ylabel={$y$}]

\addplot[green!60!black, samples=400, thick] {1 + 4/(x^2+1)};
\addlegendentry{$f_1=1+\dfrac{4}{x^2+1}$}

\addlegendimage{empty legend}\addlegendentry{}

\addplot[blue, samples=400, thick] {x^3};
\addlegendentry{$f_2=x^3$}

\addplot[red, samples=400, thick] {2^(-x)};
\addlegendentry{$f_3=2^{-x}$}

% Закрашиваем фигуру
\addplot[blue!15, samples=400, name path=F1a] {1 + 4/(x^2+1)};
\addplot[red,     samples=400, name path=F3a,  domain=-1.31:0.827] {2^(-x)};
\addplot[blue,    samples=400, name path=F2b,  domain=0.827:1.344] {x^3};

\addplot[blue!25] fill between[of=F1a and F3a, soft clip={domain=-1.31:0.827}];
\addplot[blue!25] fill between[of=F1a and F2b, soft clip={domain=0.827:1.344}];

% Вертикальные пунктиры у вершин
\addplot[dashed, thin] coordinates {(-1.31,0) (-1.31,2.47)};
\addplot[dashed, thin] coordinates {(0.827,0) (0.827,0.562)};
\addplot[dashed, thin] coordinates {(1.344,0) (1.344,2.43)};

\end{axis}
\end{tikzpicture}
\caption{Плоская фигура, ограниченная графиками $f_1$, $f_2$, $f_3$}
\label{plot1}
\end{figure}

\subsection{Анализ функций и выбор отрезков}

Функция $f_1(x)=1+4/(x^2+1)$ чётная, убывающая от $5$ (при $x=0$) до $1$
при $|x|\to\infty$.
Функция $f_2(x)=x^3$ строго возрастает: при $x<0$ отрицательна,
при $x>0$ положительна.
Функция $f_3(x)=2^{-x}$ строго убывает: при $x<0$ превышает~$1$,
при $x>0$ меньше~$1$.

Из анализа значений устанавливаем знаки разностей на концах каждого
поискового отрезка.

\medskip
\textbf{Пара $f_1\cap f_3$ (отрезок $[-2,\,-1]$):}\\
$F_{13}(x)=f_1(x)-f_3(x)$.
\begin{itemize}
    \item $F_{13}(-2)=1{,}8-4=-2{,}2<0$.
    \item $F_{13}(-1)=3-2=1>0$.
\end{itemize}
Условие знакопеременности выполнено.

\textbf{Пара $f_2\cap f_3$ (отрезок $[0,\,1]$):}\\
$F_{23}(x)=f_2(x)-f_3(x)$.
\begin{itemize}
    \item $F_{23}(0)=0-1=-1<0$.
    \item $F_{23}(1)=1-0{,}5=0{,}5>0$.
\end{itemize}

\textbf{Пара $f_1\cap f_2$ (отрезок $[1,\,2]$):}\\
$F_{12}(x)=f_1(x)-f_2(x)$.
\begin{itemize}
    \item $F_{12}(1)=3-1=2>0$.
    \item $F_{12}(2)=1{,}8-8=-6{,}2<0$.
\end{itemize}

На каждом отрезке функция $F$ монотонна (это проверяется по знаку
производной $F'$ на соответствующем отрезке; см. ниже производные),
поэтому каждый корень единственен и применимы оба выбранных метода.

\subsection{Производные функций}

\[
f_1'(x)=-\frac{8x}{(x^2+1)^2},\qquad
f_2'(x)=3x^2,\qquad
f_3'(x)=-\ln 2\cdot 2^{-x}.
\]

Монотонность $F_{13}'$ на $[-2,-1]$:
$F_{13}'(x)=f_1'(x)-f_3'(x)=-8x/(x^2+1)^2+\ln 2\cdot 2^{-x}$.
При $x\in[-2,-1]$: первое слагаемое $>0$ (так как $x<0$),
второе $>0$ (так как $2^{-x}>0$), поэтому $F_{13}'>0$. $\checkmark$

Монотонность $F_{23}'$ на $[0,1]$:
$F_{23}'(x)=3x^2+\ln 2\cdot 2^{-x}>0$. $\checkmark$

Монотонность $F_{12}'$ на $[1,2]$:
$F_{12}'(x)=-8x/(x^2+1)^2-3x^2<0$. $\checkmark$

\subsection{Выбор $\varepsilon_1$ и $\varepsilon_2$}

Площадь $S$ выражается через три корня $x_{13},x_{23},x_{12}$
и суммой двух интегралов.
Погрешность площади оценивается через погрешности корней и погрешности
интегралов~\cite{math}.
Поскольку функции ограничены и не имеют особенностей на рассматриваемых
отрезках, достаточно взять $\varepsilon_1=\varepsilon_2=\varepsilon$.
В программе \texttt{root()} и \texttt{integral()} вызываются с
одинаковым значением~$\varepsilon$.

\subsection{Разбиение площади на интегралы}

Обозначим найденные корни:
\[
    a = x_{13},\quad b = x_{23},\quad c = x_{12},
\]
где $a < b < c$.

На $[a,b]$: $f_1>f_3>f_2$ (верхняя граница — $f_1$, нижняя — $f_3$).\\
На $[b,c]$: $f_1>f_2>f_3$ (верхняя граница — $f_1$, нижняя — $f_2$).

Следовательно:
\[
    S = \int_a^b \bigl(f_1(x)-f_3(x)\bigr)\,dx
      + \int_b^c \bigl(f_1(x)-f_2(x)\bigr)\,dx.
\]

\newpage

% ================================================================== 3
\section{Результаты экспериментов}

\subsection{Точки пересечения кривых}

В таблице~\ref{table1} приведены координаты вершин фигуры,
вычисленные программой при $\varepsilon=10^{-9}$.

\begin{table}[h]
\centering
\begin{tabular}{|c|c|c|}
\hline
Кривые & $x$ & $y$ \\
\hline
$f_1$ и $f_3$ & $-1{,}306\,339\,767$ & $2{,}469\,..$ \\
$f_2$ и $f_3$ & $\phantom{-}0{,}826\,678\,573$ & $0{,}562\,..$ \\
$f_1$ и $f_2$ & $\phantom{-}1{,}344\,371\,908$ & $2{,}429\,..$ \\
\hline
\end{tabular}
\caption{Координаты точек пересечения (вывод \texttt{./a.out --roots})}
\label{table1}
\end{table}

\subsection{Площадь фигуры}

\begin{table}[h]
\centering
\begin{tabular}{|c|c|}
\hline
$\varepsilon$ & $S$ \\
\hline
$10^{-5}$ & $6{,}578\,2\ldots$ \\
$10^{-7}$ & $6{,}578\,2\ldots$ \\
$10^{-9}$ & $6{,}578\,2\ldots$ \\
\hline
\end{tabular}
\caption{Площадь фигуры при различных $\varepsilon$ (вывод \texttt{./a.out})}
\label{table2}
\end{table}

Точные значения (с 15 знаками) выводит программа при запуске без аргументов.

\newpage

% ================================================================== 4
\section{Структура программы и спецификация функций}

\subsection{Модули}

\begin{itemize}
    \item \texttt{asm/funcs.asm} — реализация на NASM32/x87
          функций $f_1,f_2,f_3,f_1',f_2',f_3'$.
    \item \texttt{src/root.c} — функция \texttt{root()}; два метода,
          выбираемые препроцессором.
    \item \texttt{src/integral.c} — функция \texttt{integral()};
          метод Симпсона с правилом Рунге.
    \item \texttt{src/area.c} — функции \texttt{compute\_area()} и
          \texttt{compute\_roots()}: использует \texttt{root()} и
          \texttt{integral()} для вычисления площади.
    \item \texttt{src/main.c} — точка входа; разбор аргументов CLI,
          тесты.
    \item \texttt{include/funcs.h}, \texttt{root.h},
          \texttt{integral.h}, \texttt{area.h} — заголовочные файлы.
\end{itemize}

\subsection{Ключевые функции}

\begin{lstlisting}
/* Нахождение корня: f(x) = g(x) на [a,b] с точностью eps */
long double root(
    long double (*f)(long double),
    long double (*g)(long double),
    long double (*df)(long double),
    long double (*dg)(long double),
    long double a, long double b, long double eps);

/* Определённый интеграл: simpson + Runge */
long double integral(
    long double (*f)(long double),
    long double a, long double b, long double eps);

/* Площадь фигуры */
long double compute_area(long double eps);
\end{lstlisting}

Число итераций после каждого вызова доступно через глобальные
переменные \texttt{root\_iters} и \texttt{integral\_iters}.

\newpage

% ================================================================== 5
\section{Сборка программы (Make-файл)}

Зависимости модулей:

\begin{center}
\texttt{a.out} $\leftarrow$ \{
\texttt{main.o}, \texttt{root.o}, \texttt{integral.o},
\texttt{area.o}, \texttt{funcs.o}
\}
\end{center}

\begin{lstlisting}[language=make,basicstyle=\ttfamily\small]
CC        = gcc
NASM      = nasm
CFLAGS    = -m32 -Wall -Wextra -Iinclude -lm
NASMFLAGS = -f elf32
METHOD   ?= NEWTON

SRC = src/main.o src/root.o src/integral.o src/area.o
ASM = asm/funcs.o

all: a.out

a.out: $(SRC) $(ASM)
	$(CC) $(CFLAGS) -DMETHOD_$(METHOD) $^ -o $@

asm/funcs.o: asm/funcs.asm
	$(NASM) $(NASMFLAGS) $< -o $@

%.o: %.c
	$(CC) $(CFLAGS) -DMETHOD_$(METHOD) -c $< -o $@

clean:
	rm -f src/*.o asm/*.o a.out
\end{lstlisting}

Команды сборки:
\begin{lstlisting}[language=bash]
make METHOD=NEWTON    # метод Ньютона
make METHOD=COMBINED  # комбинированный метод
make clean
\end{lstlisting}

\newpage

% ================================================================== 6
\section{Описание методов нахождения корней}

\subsection{Метод Ньютона (касательных)}

Пусть $F(x)=f(x)-g(x)$, $F'(x)=f'(x)-g'(x)$.
Итерационная формула~\cite{math}:
\[
    x_{n+1} = x_n - \frac{F(x_n)}{F'(x_n)}.
\]
Начальное приближение $x_0$ выбирается из условия $F(x_0)\cdot F''(x_0)>0$,
которое гарантирует сходимость и монотонное приближение к корню
с одной стороны~\cite{math}.
Второй производная $F''$ вычисляется численно по разностной схеме:
\[
    F''(x)\approx\frac{F(x+h)-2F(x)+F(x-h)}{h^2},\quad h=10^{-5}.
\]
Критерий остановки: $|x_{n+1}-x_n|<\varepsilon$.

\subsection{Комбинированный метод (хорды + касательные)}

Поддерживаются два приближения $x_N$ (сторона Ньютона)
и $x_C$ (сторона хорды). На каждой итерации:
\[
    x_N \leftarrow x_N - \frac{F(x_N)}{F'(x_N)}, \qquad
    x_C \leftarrow x_C - F(x_C)\cdot\frac{x_N-x_C}{F(x_N)-F(x_C)}.
\]
Сторона Ньютона~--- та, где $F(x_0)\cdot F''(x_0)>0$.
Метод сходится квадратично со стороны $x_N$~\cite{math}.
Критерий остановки: $|x_N-x_C|<\varepsilon$.

\newpage

% ================================================================== 7
\section{Описание метода интегрирования}

\subsection{Метод Симпсона}

Для $n$ равных разбиений отрезка $[a,b]$ (шаг $h=(b-a)/n$):
\[
    I_n = \frac{h}{3}\Bigl(f(a)+4f(a{+}h)+2f(a{+}2h)+\cdots+4f(b{-}h)+f(b)\Bigr).
\]
Погрешность метода Симпсона~\cite{math}: $|I-I_n|=O(h^4)$.

\subsection{Правило Рунге}

Начиная с $n=4$, удваиваем число разбиений до тех пор, пока
\[
    \frac{|I_{2n}-I_n|}{15} < \varepsilon.
\]
Делитель $15=2^4-1$ соответствует порядку точности $p=4$
формулы Симпсона.

\newpage

% ================================================================== 8
\section{Отладка и тестирование}

\subsection{Тесты функции \texttt{root()}}

Для каждого теста существует аналитическое значение корня.

\textbf{Тест 1} (нетривиальный): $\sin x = 0$ на $[3;\,4]$.
Аналитически: $x^* = \pi \approx 3{,}141\,592\,653\,589\,793$.
Производная: $\cos x$; второй производная: $-\sin x < 0$ на $[3,4]$,
значит начинаем с $x_0 = 3$ (там $\sin(3)>0$, $-\sin(3)<0$, произведение
$<0$) или с $x_0=4$ (там $\sin(4)<0$, $-\sin(4)>0$, произведение
$<0$)~--- начальная точка выбирается автоматически.

\textbf{Тест 2} (нетривиальный): $x^2 = 2$ на $[1;\,2]$.
Аналитически: $x^* = \sqrt{2}\approx 1{,}414\,213\,562\,373\,095$.
$F(x)=x^2-2$, $F'(x)=2x$, $F''(x)=2>0$;
$F(1)=-1<0$, $F(2)=2>0$;
$F(2)\cdot F''(2)=2\cdot2=4>0 \Rightarrow x_0=2$.

\textbf{Тест 3} (нетривиальный): $e^x = 3$ на $[0;\,2]$.
Аналитически: $x^* = \ln 3\approx 1{,}098\,612\,288\,668\,110$.
$F'(x)=e^x>0$, $F''(x)=e^x>0$;
$F(0)=-2<0$, $F(2)=e^2-3>0$;
$F(2)\cdot F''(2)>0 \Rightarrow x_0=2$.

Результаты при $\varepsilon=10^{-9}$ (вывод \texttt{./a.out --test-root}):

\begin{table}[h]
\centering
\begin{tabular}{|l|c|c|c|}
\hline
Тест & Результат & Аналитически & $|\text{ошибка}|$ \\
\hline
$\sin x=0,\,[3,4]$  & $3{,}141592653\ldots$ & $\pi$ & $<10^{-12}$ \\
$x^2=2,\,[1,2]$     & $1{,}414213562\ldots$ & $\sqrt{2}$ & $<10^{-12}$ \\
$e^x=3,\,[0,2]$     & $1{,}098612288\ldots$ & $\ln 3$ & $<10^{-12}$ \\
\hline
\end{tabular}
\caption{Результаты тестов \texttt{root()}}
\label{table3}
\end{table}

\subsection{Тесты функции \texttt{integral()}}

\textbf{Тест 1}: $\displaystyle\int_0^1 x^2\,dx = \frac{1}{3}$.

\textbf{Тест 2} (нетривиальный): $\displaystyle\int_0^\pi \sin x\,dx = 2$.
(Метод Симпсона точен для полиномов степени $\le 3$,
поэтому тест с $\sin$ нетривиален.)

\textbf{Тест 3}: $\displaystyle\int_0^1 x^3\,dx = \frac{1}{4}$.

\begin{table}[h]
\centering
\begin{tabular}{|l|c|c|c|}
\hline
Тест & Результат & Аналитически & $|\text{ошибка}|$ \\
\hline
$\int_0^1 x^2\,dx$     & $0{,}333\,333\,333\ldots$ & $1/3$   & $<10^{-14}$ \\
$\int_0^\pi \sin x\,dx$ & $2{,}000\,000\,000\ldots$ & $2$     & $<10^{-12}$ \\
$\int_0^1 x^3\,dx$     & $0{,}250\,000\,000\ldots$ & $1/4$   & $<10^{-14}$ \\
\hline
\end{tabular}
\caption{Результаты тестов \texttt{integral()}}
\label{table4}
\end{table}

\newpage

% ================================================================== 9
\section{Программа на Си и на ассемблере}

Исходные тексты программы прилагаются к отчёту в архиве
\texttt{Moskvin.tar.gz} в составе:

\begin{itemize}
    \item \texttt{Makefile}
    \item \texttt{asm/funcs.asm}
    \item \texttt{src/main.c}, \texttt{src/root.c},
          \texttt{src/integral.c}, \texttt{src/area.c}
    \item \texttt{include/funcs.h}, \texttt{include/root.h},
          \texttt{include/integral.h}, \texttt{include/area.h}
\end{itemize}

\newpage

% ================================================================== 10
\section{Анализ допущенных ошибок}

В процессе работы были допущены и исправлены следующие ошибки.

\begin{itemize}
    \item При реализации \texttt{f3} и \texttt{f3d} на ассемблере
          необходимо учесть, что инструкция \texttt{f2xm1} требует
          аргумент в диапазоне $[-1,1]$. Нарушение этого условия
          давало неверный результат; исправлено предварительным
          разбиением показателя на целую и дробную части
          с помощью \texttt{frndint}.
    \item При использовании \texttt{fdivp} и \texttt{fdivrp}
          легко перепутать порядок делителя и делимого. Для ясности
          используется явный \texttt{fxch} перед \texttt{fdivp}.
    \item Для гарантированного выбора начальной точки в методе
          Ньютона необходим знак $F''(x_0)$; поскольку передача
          второй производной в API усложнила бы интерфейс,
          $F''$ аппроксимируется численно.
\end{itemize}

\newpage

\begin{raggedright}
\addcontentsline{toc}{section}{Список цитируемой литературы}
\begin{thebibliography}{9}
\bibitem{math}
    Ильин~В.\,А., Садовничий~В.\,А., Сендов~Бл.\,Х.
    Математический анализ. Т.\,1.~---
    Москва: Наука, 1985.
\end{thebibliography}
\end{raggedright}

\end{document}
